% Part: first-order-logic
% Chapter: axiomatic-deduction
% Section: rules-and-proofs

\documentclass[../../../include/open-logic-section]{subfiles}

\begin{document}

\iftag{FOL}
      {\olfileid[fr]{fol}{axd}{rul}}
      {\olfileid[fr]{pl}{axd}{rul}}

\olsection{Règles et \usetoken{p}{derivation}}

\begin{explain}
Les !!{derivation}s axiomatiques constituent peut-être le système de
!!{derivation} le plus simple pour la logique. !!^a{derivation} n'est
qu'une suite de !!{formula}s. Pour être une !!{derivation}, chaque
!!{formula} de la suite doit soit être une instance d'un axiome, soit
découler, par une règle d'inférence, d'une ou de plusieurs !!{formula}s
qui la précèdent dans la suite. La dernière !!{formula}
d'!!a{derivation} est celle qu'elle permet de !!{derive}.
\end{explain}

\begin{defn}[!!^{derivability}]
Si $\Gamma$ est un ensemble de !!{formula}s de $\Lang L$, une
\emph{!!{derivation} à partir de~$\Gamma$} est une suite finie
$!A_1$, \dots,~$!A_n$ de !!{formula}s telle que, pour chaque
$i \le n$, l'une des conditions suivantes soit satisfaite :
\begin{enumerate}
\item $!A_i \in \Gamma$;
\item $!A_i$ est un axiome;
\item $!A_i$ découle, par une règle d'inférence, d'une formule~$!A_j$
  (et d'une formule~$!A_k$) telles que $j < i$ (et $k < i$).
\end{enumerate}
\end{defn}

Ce qui constitue une !!{derivation} correcte dépend des règles
d'inférence que nous admettons, ainsi, bien entendu, que des axiomes
que nous choisissons. Une règle d'inférence est un énoncé conditionnel
qui indique sous quelles conditions une étape~$!A_i$ d'!!a{derivation}
est une étape d'inférence correcte.

\begin{defn}[Règle d'inférence]
Une \emph{règle d'inférence} donne une condition suffisante pour
qu'une étape d'!!a{derivation} à partir de~$\Gamma$ soit une étape
d'inférence correcte.
\end{defn}

Par exemple, toute suite à un seul terme~$!A$ telle que
$!A \in \Gamma$ est trivialement une !!{derivation}. On pourrait donc
formuler la règle d'inférence très simple suivante :
\begin{quote}
  Si $!A \in \Gamma$, alors $!A$ est toujours une étape d'inférence
  correcte dans toute !!{derivation} à partir de~$\Gamma$.
\end{quote}
De même, si $!A$ est l'un des axiomes, alors $!A$ constitue à lui seul
une !!{derivation}; ceci donne également une règle d'inférence :
\begin{quote}
  Si $!A$ est un axiome, alors $!A$ est une étape d'inférence correcte.
\end{quote}
Les règles qui font appel à des !!{formula}s antérieures à l'étape
considérée sont plus intéressantes. La règle suivante s'appelle le
\emph{modus ponens} :
\begin{quote}
  Si $!B \lif !A$ et $!B$ figurent plus haut dans la !!{derivation},
  alors~$!A$ est une étape d'inférence correcte.
\end{quote}
Si le modus ponens est la seule règle d'inférence, la définition
précédente revient à dire que $!A_1$, \dots,~$!A_n$ est une
!!{derivation} si et seulement si, pour chaque $i \le n$, l'une des
conditions suivantes est satisfaite :
\begin{enumerate}
\item $!A_i \in \Gamma$;
\item $!A_i$ est un axiome;
\item pour un certain $j < i$, $!A_j$ est $!B \lif !A_i$, et, pour un
  certain $k < i$, $!A_k$ est~$!B$.
\end{enumerate}
La dernière clause affirme que $!A_i$ découle de~$!A_j$
($!B \lif !A_i$) et de~$!A_k$ ($!B$) par modus ponens. Si nous pouvons
parcourir les indices de~$1$ à~$n$ et constater, à chaque étape, que
la !!{formula}~$!A_i$ appartient à~$\Gamma$, est un axiome, ou satisfait
une règle d'inférence, alors la suite entière est une !!{derivation}
correcte.

\begin{defn}[!!^{derivability}]
Une !!{formula}~$!A$ est \emph{!!{derivable} à partir de~$\Gamma$}, ce
que l'on note $\Gamma \Proves !A$, s'il existe une !!{derivation} à
partir de~$\Gamma$ qui se termine par~$!A$.
\end{defn}

\begin{defn}[Théorèmes]
!!^a{formula}~$!A$ est un \emph{théorème} s'il existe une
!!{derivation} de~$!A$ à partir de l'ensemble vide. On écrit
$\Proves !A$ si $!A$ est un théorème, et $\Proves/ !A$ dans le cas
contraire.
\end{defn}

\end{document}
